\section{\label{apx:dpsgd}DP-SGD implementation}
This appendix covers the accountant validation,
the precision and lot-size choices,
the per-example gradient implementation,
and the pitfalls specific to the private branch.

\paragraph{Accountant validation.}
For target $\varepsilon = 8$ at $q = 0.01$,
$100$ steps and $\delta = 10^{-5}$,
the calibration returns $\sigma \approx 0.50$ and the recomputed guarantee is $\varepsilon \approx 7.995$.
At $3$ epochs and $B = 16$ over $16{,}384$ rows ($3072$ steps, $q \approx 9.77 \times 10^{-4}$) the calibration gives $\sigma \approx 0.437$ at $\varepsilon = 8$ and $\sigma \approx 0.864$ at $\varepsilon = 1$. Re-running the accounting with R\'enyi orders up to $1024$ reproduces the realized $\varepsilon$ to three decimals at both budgets,
so the default order grid costs no utility.

\paragraph{Numerical precision.}
The lot gradient is accumulated into an fp32 buffer and cast back to bf16 only at gradient assignment.
Performing the accumulation and the noise draw in bf16 instead rounds much of the calibrated noise away,
and the run then reports an $\varepsilon$ it does not deliver.

\paragraph{Lot size.}
At a lot size of $16$ the private conditions are severely under-batched.
With $\sigma = 0.437$ and $d = 3{,}211{,}264$ trainable parameters,
the averaged noise vector has norm $\approx 49$ against a signal of at most $C = 1$:
the update is roughly $98\%$ noise.
Lot size is the only knob that fixes this for free.
Analytically,
$\sigma$ grows only about $7\%$ per doubling of $B$,
so $B = 256$ at $3$ epochs yields $12.1\times$ the gradient signal-to-noise ratio at $1.0\times$ the cost,
whereas $B = 256$ at $6$ epochs yields \textit{less} ($11.3\times$) at $2.0\times$ the cost.
The measured ratio improved about $7\times$ ($0.010 \to 0.073$ at $\varepsilon = 8$).
Raising the adapter rank instead would have been actively harmful,
since the noise is injected per trainable parameter and more capacity therefore means more noise at the same $\varepsilon$.

\paragraph{Per-example gradients by explicit loop.}
The natural implementation vectorizes per-example gradients with \texttt{torch.func}~\citep{paszke2019pytorch},
and the small adapter parameter count makes that look cheap.
It is not.
There is no batching rule for the fused FlashAttention~\citep{dao2022flashattention} backward kernel,
so the vectorizing transform falls back to a path that materializes the full sequence\,$\times$\,sequence attention matrix per example,
and the branch exhausted memory at $16 \times 2048$ on the grid instance.
Neither a lower LoRA rank nor micro-batching addresses this,
because the vectorization itself is the cost.
The private branch therefore iterates over the lot one example at a time,
clipping in place and accumulating.
This is slower per step but bounded,
and the ratio is structural rather than tunable -- one backward pass per example regardless of lot size -- costing about $2.8\times$ the wall-clock of the non-private mechanism ($97$ against $35$ minutes at $B = 16$).

\paragraph{Lot size variability.}
Under Poisson subsampling the realized lot size varies (at $B = 16$: mean $16$, sd $3.9$, observed range $7$--$30$),
and empty lots are skipped,
which only over-pays privacy.

\paragraph{Two pitfalls specific to the private branch.}
Rows whose prompt fills the context window have every label masked,
so their per-example cross-entropy is NaN;
summed into the lot gradient this contaminates every adapter parameter and yields a garbage checkpoint with only a NaN in the loss log to show for it.
Such rows are dropped before training and counted.
Separately,
dropout is disabled by setting every dropout module's rate to zero rather than by placing the model in evaluation mode:
a live dropout mask would break the per-example determinism the accounting assumes,
while evaluation mode would additionally disable training-mode behavior elsewhere in the network.

